\documentclass[preview]{standalone}
\usepackage{amsmath}
\usepackage{booktabs}
\usepackage{amssymb}
\usepackage{setspace}
\usepackage{bm}
\let\oldcdot\cdot
\usepackage{breqn}
\let\cdot\oldcdot
\begin{document}
Let $\phi(x)$ be the level-set function and
 \begin{align*}\raisebox{\depth}{$\chi$}(x) &= \begin{cases} 1 & \mbox{if}\quad \phi(x) < 0 \\ k_{air}/k_0 & \mbox{if}\quad \phi(x)  \ge 0 \end{cases},\end{align*}
where the weak-material contrast is the physically motivated ratio of the thermal conductivity of air to that of aluminium, $k_{air}/k_0 = 0.026/200 \approx 1.3 \times 10^{-4}$ (numerically $1.3 \times 10^{-4}$, the value taken by the characteristic function).
The design domain is the normalized unit square — the cross-section of the conducting wall of a flue-gas heat-recovery duct scaled to unit size for benchmark comparison.  Low-temperature exhaust gas at $u_H = 338$ K sweeps the central forty percent of the bottom edge and injects heat by TWO parallel mechanisms: convection through a Robin condition with heat transfer coefficient $h_H = 25$ W/(m$^{2}$ K) AND gas radiation — the wall ABSORBS the radiative flux $q_{rad} = \varepsilon_g \sigma \left( u_{H}^{4} - u^{4}\right)$ emitted by the hot gas, where $\varepsilon_g = 0.9$ is the effective emissivity of the gas-filled duct and the oxidised wall and $\sigma = 5.67 \times 10^{-8}$ W/(m$^{2}$ K$^{4}$) the Stefan--Boltzmann constant; when the wall is cold the radiative flux reaches $\varepsilon_g \sigma u_H^{4} \approx 666$ W/m$^{2}$, of the same order as the convective load $h_H\left(u_H - u_{amb}\right) = 1000$ W/m$^{2}$.  The remaining bottom edge is adiabatic and the bottom edge is a fixed boundary (the gas-side footprint is prescribed).  The other three edges exchange heat with ambient air at $u_{amb} = 298$ K through an effective convective heat transfer coefficient $h = 25$ W/(m$^{2}$ K), representative of the duct casing in ambient air.  The $u^{4}$ dependence of the absorbed flux makes the state problem NONLINEAR and is the physical mechanism studied here.  The optimisation minimises the gas-side surface temperature at fixed material cost.
The Lagrangian of the flue-gas problem is
\begin{dmath*}
\mathcal L(\left[ u\right],\left[ \delta u\right]) = J(u) +  \int_{D} k \raisebox{\depth}{$\chi$} \left(  \nabla u  \cdot  \nabla \delta u  \right) \,dx + \int_{\Gamma_H} h_H \left( u - u_{H}\right) \delta u \,ds_x - \int_{\Gamma_H} \varepsilon_g \sigma \left( u_{H}^{4} - u^{4}\right) \delta u \,ds_x + \int_{\Gamma_c} h \left( u - u_{amb}\right) \delta u \,ds_x ,
\end{dmath*}
with
\begin{align*}
J(u) &=  \int_{D} \ell \raisebox{\depth}{$\chi$} \,dx + \int_{\Gamma_H} q_b u \,ds_x ,
 \end{align*}
where the heated-surface temperature integral $\int_{\Gamma_H} u \,ds_x$ is weighted by $q_b = h_H \left( u_H - u_{amb}\right)$, the heat flux the heated side would absorb if the surface reached ambient temperature; minimizing $J$ keeps the heated surface cold at fixed material cost, i.e. it maximizes the heat absorbed from the hot fluid.  The Lagrange multiplier $\ell$ weights the material cost; it is a balance parameter of the optimisation, calibrated by the pipeline (a geometric scan locating the knee of the volume-fraction curve), not part of the problem statement.  There is no Dirichlet condition: the temperature is fixed implicitly by the ambient air through the Robin condition.
The state $\left[ u\right] \in \mathbb P_{2}(D)$ is the solution of the NON-LINEAR problem
\begin{dmath*}
\int_{D} k \raisebox{\depth}{$\chi$} \left(  \nabla u  \cdot  \nabla \delta u  \right) \,dx + \int_{\Gamma_H} h_H u \delta u \,ds_x + \int_{\Gamma_H} \varepsilon_g \sigma u^{4} \delta u \,ds_x + \int_{\Gamma_c} h u \delta u \,ds_x = \int_{\Gamma_H} h_H u_{H} \delta u \,ds_x + \int_{\Gamma_H} \varepsilon_g \sigma u_{H}^{4} \delta u \,ds_x + \int_{\Gamma_c} h u_{amb} \delta u \,ds_x \qquad \forall\; {\delta u\in \mathbb P_{2}(D)} .
\end{dmath*}
The problem is solved by Newton's method: the state $u$ is updated as $u \leftarrow u - \hat u$, where the Newton update $\left[ \hat u\right]$ at the state $\left[ u\right]$ is the solution of
\begin{dmath*}
\int_{D} k \raisebox{\depth}{$\chi$} \left(  \nabla \hat u  \cdot  \nabla \delta u  \right) \,dx + \int_{\Gamma_H} h_H \hat u \delta u \,ds_x + \int_{\Gamma_H} 4 \varepsilon_g \sigma u^{3} \hat u \delta u \,ds_x + \int_{\Gamma_c} h \hat u \delta u \,ds_x = \int_{D} k \raisebox{\depth}{$\chi$} \left(  \nabla u  \cdot  \nabla \delta u  \right) \,dx + \int_{\Gamma_H} \left( h_H \left( u - u_H\right) + \varepsilon_g \sigma \left( u^{4} - u_{H}^{4}\right)\right) \delta u \,ds_x + \int_{\Gamma_c} h \left( u - u_{amb}\right) \delta u \,ds_x \qquad \forall\; {\delta u\in \mathbb P_{2}(D)} ,
\end{dmath*}
where the radiation term is linearised as $\partial_{u} \left[ \varepsilon_g \sigma \left( u^{4} - u_{H}^{4}\right)\right] \hat u = 4 \varepsilon_g \sigma u^{3} \hat u$.
The adjoint state $\left[ \tilde u\right]$ to the direct state $\left[ u\right]$ is the solution of
\begin{dmath*}
\int_{D} k \raisebox{\depth}{$\chi$} \left(  \nabla \delta u  \cdot  \nabla \tilde u  \right) \,dx + \int_{\Gamma_H} h_H \delta u \tilde u \,ds_x + \int_{\Gamma_H} 4 \varepsilon_g \sigma u^{3} \delta u \tilde u \,ds_x + \int_{\Gamma_c} h \delta u \tilde u \,ds_x = - \int_{\Gamma_H} q_b \delta u \,ds_x \qquad \forall\; {\delta u\in \mathbb P_{2}(D)} ,
\end{dmath*}
where the adjoint operator is the linearisation of the state operator at $u$ (self-adjoint for the symmetric conduction--convection--radiation operator), and the adjoint load $- \int_{\Gamma_H} q_b \delta u \,ds_x$ comes from the partial derivative $\partial_{u} J(\delta u) = \int_{\Gamma_H} q_b \delta u \,ds_x$ of the heated-surface temperature integral in the objective.
In order to compute the shape derivative, we consider a shape perturbation $\mathbf x = T_t(\mathbf X) = \mathbf X + t \mathbf  V(\mathbf X)$ with a suitable velocity field $\mathbf V$, vanishing on the fixed gas-side footprint.  The perturbed Lagrangian $\mathcal G$ using the Lagrangian state $\left[ u\right]^{L}$ is given by
\begin{dmath*}
\mathcal G(t,\left[ u\right]^{L},\left[ \delta u\right]^{L}) = \int_{D} \ell \raisebox{\depth}{$\chi$}  \operatorname{det}\mathbf F( t )  \,dx + \int_{\Gamma_H} q_b u^{L} \,ds_x +  \int_{D} k \raisebox{\depth}{$\chi$}  \operatorname{det}\mathbf F( t )  \left( \left(  \nabla u^{L} \cdot  \mathbf F^{-1}(t) \right) \cdot \left(  \nabla \delta u^{L} \cdot  \mathbf F^{-1}(t) \right) \right) \,dx + \int_{\Gamma_H} h_H \left( u^{L} - u_{H}\right) \delta u^{L} \,ds_x - \int_{\Gamma_H} \varepsilon_g \sigma \left( u_{H}^{4} - \left( u^{L}\right)^{4}\right) \delta u^{L} \,ds_x + \int_{\Gamma_c} h \left( u^{L} - u_{amb}\right) \delta u^{L} \,J_{s}(t) \,ds_x ,
\end{dmath*}
where $\mathbf F(t) = \mathbf I + t \nabla \mathbf V$ is the deformation gradient and $J_{s}(t)$ the boundary Jacobian of the moving cooling edges $\Gamma_c$; the fixed gas side $\Gamma_H$ carries no Jacobian because $\mathbf V$ vanishes there, and BOTH the convection and the radiation integrals over $\Gamma_H$ therefore contribute NOTHING to the shape derivative.
For the direct state $\left[ u\right]$ and the adjoint state $\left[ \tilde u\right]$, the volume expression of the shape derivative, including the curvature term on the moving cooling edges, is given by
\begin{dmath*}
DJ(\Omega)(\mathbf V) = \int_{D} \left(\operatorname{div} {\mathbf V} \left(  \nabla u  \cdot  \nabla \tilde u  \right) - \left(  \nabla \tilde u  \cdot \left(  \nabla u  \cdot  \nabla {\mathbf V}  \right) \right) - \left(  \nabla u  \cdot \left(  \nabla \tilde u  \cdot  \nabla {\mathbf V}  \right) \right)\right) k \raisebox{\depth}{$\chi$} \,dx + \int_{D} \ell \raisebox{\depth}{$\chi$} \operatorname{div} {\mathbf V} \,dx + \int_{\Gamma_c} \kappa h \left( u - u_{amb}\right) \tilde u \left(\mathbf V \cdot \mathbf n\right) \,ds_x ,
\end{dmath*}
where $\kappa$ is the curvature of $\Gamma_c$ and $\mathbf n$ its outward unit normal.  The gas-side footprint is a fixed boundary; the cooling edges $\Gamma_c$ are design boundaries — their shape is optimised.
The initial design is the full solid square with a circular hole of radius 0.1 cut out of its centre.  The temperature field uses P2 finite elements; there is no Dirichlet condition, the temperature is fixed implicitly by the ambient air through the Robin condition.
\end{document}
